\chapter{Verified Rewrite-Based Optimization and Instruction Selection}
\label{ch:rewrite-optimization}

An algebraic identity and a profitable instruction pattern are different claims.
The identity says two computations have the same permitted behavior.
The pattern says a target sequence implements one computation at a chosen cost.
A rewrite-based compiler needs both claims, plus a selection algorithm whose optimality domain is stated honestly.

\section{Bind rules to an arithmetic and effect model}
\label{sec:ch88-semantics}
In total mathematical integer arithmetic, add(x,zero)$\to$x and mul(x,one)$\to$x preserve the value and still evaluate x once.
Multiplying x by zero and replacing the whole expression by zero can remove evaluation of x.
If x diverges or performs an observable action, that change is invalid even when arithmetic is exact.
The functional optimizer therefore deliberately omits that rule in its strict language.

Floating-point rules need a different analysis.
Reassociation changes rounding, signed zeros can distinguish some algebraic forms, and NaNs affect comparison identities.
Checked integer arithmetic introduces overflow behavior that modular arithmetic does not have.
Table~\ref{tab:ch88-rules} lists the rule's semantic premises rather than pretending a textual pattern implies them.

\begin{table}[htbp]
\caption{A rewrite is valid only under the computation contract it preserves.}
\label{tab:ch88-rules}
\centering
\begin{tabular}{ll}
\toprule
Candidate & Required contract \\
\midrule
add(x,zero) to x & compatible arithmetic; x still evaluated as required \\
mul(x,one) to x & compatible arithmetic; x still evaluated as required \\
mul(x,zero) to zero & discarded computation has no required behavior \\
reassociate additions & arithmetic permits changed grouping \\
fuse multiply-add & target rounding/overflow semantics match source \\
reuse load & storage and intervening effects justify reuse \\
\bottomrule
\end{tabular}
\end{table}

A target fused multiply-add instruction is not automatically equivalent to separate floating multiply and add because it can round only once.
The lab's selector is a structural cost model for an explicitly suitable pure arithmetic interpretation, not a floating-point fusion authorization.
Instruction availability, operand constraints, and exception behavior belong to the target rule specification.

\section{Validate finite rules independently of search}
\label{sec:ch88-validation}
Algorithm~\ref{alg:ch88-check} checks a proposed rule over an exhaustive finite input domain.
For $V$ variables each ranging over $D$ values, it evaluates $D^V$ assignments and compares defined results and observable traces.
This is a complete finite-domain check when the evaluator and enumeration correctly implement that domain, but not a proof for unbounded integers or arbitrary heap states.
The earlier byte-arithmetic e-graph lab provides such exhaustive finite checks for its identities.

\begin{algorithm}[htbp]
\caption{Validate a rewrite over a stated finite semantic domain.}
\label{alg:ch88-check}
\begin{pseudocode}
\Require rule, finite variable domains, independent semantics including failures/effects
\Ensure counterexample or exhaustive agreement for that domain only
\ForAll{assignments to rule variables}
\State evaluate both sides under the same input state and applicable premises
\State compare values, final observable state, errors, and termination classification
\If{outcomes differ} \State return the assignment as counterexample \EndIf
\EndFor
\State record domain and evaluator version with the verified rule
\end{pseudocode}
\end{algorithm}

For unrestricted programs, evaluating both sides need not terminate, so the finite algorithm also needs a bounded semantics or a decidable total fragment.
Treating evaluator timeout on both sides as semantic equivalence is unsound.
Solver-based validation can cover larger domains, but its undefined-behavior, poison, arithmetic, and memory encodings must match the IR.
The proof boundary remains separate from profitability search.

Figure~\ref{fig:ch88-pipeline} places validation before target ranking and after final lowering where appropriate.
A rewrite certificate can establish that the engine applied a known rule correctly while leaving the rule's semantic proof in a reusable trusted database.
That reduces the checker burden without making an incorrect trusted rule harmless.

\begin{figure}[htbp]
\centering
\begin{tikzpicture}[node distance=12mm]
\node[basic block] (sem) {semantic rule and premises};
\node[basic block,below=of sem] (proof) {proof or finite-domain validation};
\node[basic block,below=of proof] (search) {matching and cost-based selection};
\node[basic block,below=of search] (emit) {constraint repair and emitted-code checks};
\draw[cfg edge] (sem)--(proof);\draw[cfg edge] (proof)--(search);\draw[cfg edge] (search)--(emit);
\end{tikzpicture}
\caption{A favorable target cost cannot substitute for a semantic justification.}
\label{fig:ch88-pipeline}
\end{figure}

\section{Select an additive tree cover}
\label{sec:ch88-select}
The lab target assigns cost zero to available leaves, two to add, three to multiply, and three to a fused multiply-add pattern matching add(mul(a,b),c).
For that expression, separate selection costs five while the fused cover costs three.
The pattern consumes the add and multiply together; its frontier children are a, b, and c.
Table~\ref{tab:ch88-costs} works both covers.

\begin{table}[htbp]
\caption{A combined target rule can cover several source nodes at lower additive cost.}
\label{tab:ch88-costs}
\centering
\begin{tabular}{lrrr}
\toprule
Cover & Root rule & Frontier cost & Total \\
\midrule
separate multiply and add & add: 2 & multiply: 3 & 5 \\
combined multiply-add & madd: 3 & leaves: 0 & 3 \\
\bottomrule
\end{tabular}
\end{table}

Algorithm~\ref{alg:ch88-dp} minimizes the additive tree-cover cost.
With a fixed finite pattern set, memoized or bottom-up dynamic programming stores one best cost per node and result category and evaluates each matching cover once.
For $N$ nodes and $P$ fixed-size patterns, a direct implementation takes $O(NP)$ matching/cost work plus reconstruction output, assuming bounded categories and constant-time pattern tests.
The tiny executed selector recursively enumerates its few alternatives without a general memoization table, so it is a correctness fixture for these examples rather than an implementation of that production complexity bound.

\begin{algorithm}[htbp]
\caption{Compute a minimum additive target cover for a finite expression tree.}
\label{alg:ch88-dp}
\begin{pseudocode}
\Require tree, finite legal target patterns, additive nonnegative costs, result categories
\Ensure minimum-cost cover within this tree grammar
\ForAll{nodes in postorder}
\ForAll{target patterns matching this node}
\State obtain required categories and best costs at the pattern frontier
\State add frontier costs to this target rule's cost
\State retain the cheapest feasible cover for each produced category
\EndFor
\EndFor
\State reconstruct the selected root cover through recorded backpointers
\end{pseudocode}
\end{algorithm}

Figure~\ref{fig:ch88-cover} identifies the combined region.
If register classes, immediate ranges, or two-address constraints matter, categories and side conditions must express them or a later repair phase can invalidate the advertised cost.
Cost numbers in the lab are declared abstract units, not measured processor cycles.
No benchmark speedup is inferred from reducing five units to three.

\begin{figure}[htbp]
\centering
\begin{tikzpicture}[node distance=13mm and 14mm]
\node[basic block] (add) {add};\node[basic block,below left=of add] (mul) {mul};
\node[basic block,below right=of add] (c) {c};\node[basic block,below left=of mul] (a) {a};
\node[basic block,below right=of mul] (b) {b};
\draw[cfg edge] (add)--(mul);\draw[cfg edge] (add)--(c);\draw[cfg edge] (mul)--(a);\draw[cfg edge] (mul)--(b);
\node[draw,dashed,fit=(add)(mul),inner sep=7pt,label=above:{madd cover}] {};
\end{tikzpicture}
\caption{One target rule covers the add/multiply region and leaves three inputs on its frontier.}
\label{fig:ch88-cover}
\end{figure}

\section{Do not transfer tree optimality to shared DAGs}
\label{sec:ch88-dags}
Suppose a multiply result feeds two additions.
Covering it separately in two fused instructions may duplicate work; computing it once may preserve sharing but require a register and extra live range.
An additive tree model sees duplicated subtrees and can prefer a different result from a DAG-aware model.
Register allocation, scheduling, and placement make cost context-dependent even after structural covering.

A graph selector can use integer programming, bounded search, or heuristics over sharing and instruction constraints.
Any optimality claim is relative to its encoded legal alternatives and cost model, and search limits can leave the best found solution unproven optimal.
The reference lab does not implement a general DAG selector.
The existing instruction-selection and register-allocation chapters supply the larger backend setting for this distinction.

\section{Connect destructive rewriting to equality saturation}
\label{sec:ch88-egraphs}
A destructive rewrite drops an alternative unless the engine retains history.
An e-graph retains known equivalent forms and lets extraction choose later, but equality discovery and target extraction remain different phases.
Extraction from positive additive costs can find a minimum tree representative in a fixed e-graph; it cannot discover missing equalities or account automatically for shared-register cost.
Conditional e-class facts must remain valid under class merges.

The executed rewrite lab checks matching, occurs checks, certificates, overlap construction, shared updates, dangling references, and the small target cover.
The functional lab independently executes source, optimized source, and compiled stack code on many closed programs.
Together these tests target different failure modes rather than treating one successful output as evidence for every rewrite mechanism.
The expansion record states which mechanisms are implemented and which remain design discussions.

\section{Prove the dynamic program's limited optimality claim}
\label{sec:ch88-optimality}
The tree-cover recurrence is correct by induction on subtree height under its stated grammar and additive cost model.
Every legal root cover chooses one matching target rule.
That rule leaves a finite frontier of smaller subtrees with required result categories.
By induction, the stored best cover for each frontier subtree is no more expensive than any other legal cover producing that category.
Adding those costs to the root rule cost and minimizing across all root choices therefore finds the minimum admitted cover.

The proof identifies its limitations.
If the pattern index omits a legal root rule, the recurrence optimizes an incomplete grammar.
If two frontier subtrees share a computation whose cost should be paid only once, adding their independent costs overcounts it.
If a register constraint couples their selections, independently cheapest category choices may not compose into a legal machine program.
If a rule's cost depends on later scheduling, a fixed scalar cost is only an estimate.
None of these failures is repaired merely by calling the search dynamic programming.

Categories can encode some additional constraints, such as integer versus floating register classes, available immediate forms, or alternative result locations.
Increasing category precision also increases the state space.
A target description should state which constraints are encoded exactly and which are repaired heuristically after selection.
The final cost claim should refer to the selected grammar, not to all possible machine programs on the processor.

\section{Retain proof assumptions through target lowering}
\label{sec:ch88-assumptions}
A high-level equality can stop being useful when lowered operations acquire narrower semantics.
An exact integer addition may become a checked machine operation with a failure edge or a wrapping operation whose result is observed at a fixed width.
A proof established before that choice must either quantify over the intended later interpretation or be transported through a verified lowering relation.
Simply copying the same opcode spelling into a lower IR does not preserve the proof.

Memory and floating-point relaxations should be equally explicit.
A transformation admitted under an optional relaxed arithmetic flag must not leak into a strict compilation unit through an unkeyed rewrite cache.
A load-reuse rule justified by a no-alias fact must be invalidated when pointer provenance or the call graph changes.
Proof metadata should identify semantic mode, relevant facts, and source/target operation versions so a validator can reconstruct why a particular step was legal.

At the final machine boundary, differential execution is valuable but necessarily selective.
The spellchecker uses an independent distance formulation, the rewrite engine replays local certificates, and the functional compiler compares source and VM evaluation.
A native selector additionally needs target instruction semantics, register/clobber contracts, and tests at encoding boundaries.
These checks complement rather than replace one another.
The book's earlier machine-code and validation case studies provide the next layers after the abstract target-cover fixture used here.

\section{Exercises}
\label{sec:ch88-exercises}
\begin{enumerate}
\item\label{ex:ch88-cover} Compute both covers of add(mul(a,b),c) using the stated costs.
\item\label{ex:ch88-diverge} Give a strict pure program where replacing mul(e,zero) by zero changes termination.
\item\label{ex:ch88-dag} A multiply feeds two adds.
Explain why two locally cheapest fused covers need not be globally cheapest once sharing and registers matter.
\item\label{ex:ch88-proof} A rule passes all 256 byte inputs.
What additional reasoning is needed before applying it to checked 64-bit arithmetic or floating point?
\end{enumerate}

\section{Further reading}
\label{sec:ch88-reading}
Ohlebusch develops the rewrite foundations~\cite{ohlebusch-2002-rewriting}; the \code{egg} paper explains alternative-preserving search~\cite{willsey-2021-egg}.
Appel and Morgan connect tree/graph representations to target selection and optimization~\cite{appel-1998-modern,morgan-1998-optimizing}.
